\clearpage
\fancyhead[L]{\small SOLDNER. AZIMUTHE ZU REDUCIREN.}
\SourceLeaf{592}
\begin{center}\small Göttingische gelehrte Anzeigen. Stück 46. Seiten 449..456. 23. März 1815.\\[3pt]\rule{.72\linewidth}{.3pt}\end{center}

I. \emph{Neue Methode zur Reduktion beobachteter Azimute. Von} J. \textsc{Soldner}. München. 14 Seiten in Quart. (Sonderdruck aus den Denkschriften der Münchner Akademie für 1813.)

II. \emph{Beiträge zur Berechnung beobachteter Azimute von} \textsc{Anton von Stefenelli}. Ebenda. Gedruckt mit Stornoschen Lettern. 30 Seiten in Quart.

Weil sich die Azimute ungleichförmig ändern, darf man den Mittelwert einer Reihe von Azimutunterschieden zwischen einem irdischen Gegenstand und einem Himmelskörper nicht mit dem Azimutunterschied gleichsetzen, der dem Mittel aller Beobachtungszeiten entspricht. Hat man freilich nach jeder einzelnen Beobachtung abgelesen, so wird man auch die Mühe nicht scheuen, jede einzeln zu berechnen, um ebenso viele einzelne Ergebnisse zu erhalten. In der Regel liest man aber erst am Ende einer Beobachtungsreihe ab, um dafür mehr Beobachtungen anstellen zu können. Eine große Zahl von Beobachtungen einzeln zu berechnen, wäre eine abschreckend langwierige Arbeit. Würde man das Mittel aller beobachteten Azimutunterschiede mit dem Azimutunterschied zur mittleren Beobachtungszeit verwechseln, so entstünden Fehler, die man bei der Genauigkeit der Beobachtungen mit Theodoliten, wie sie \textsc{Reichenbach} gegenwärtig liefert, nicht übersehen darf. Zum Teil aus diesem Grund schlug ein berühmter praktischer Astronom vor, diese Beobachtungen nur um die Zeit der Kulmination anzustellen. Dann darf man die Änderung des Azimuts tatsächlich als gleichförmig ansehen und hat außerdem den Vorteil, dass Fehler in der Polhöhe oder der Deklination des Gestirns das Ergebnis nur ganz unmerklich beeinflussen. Bei diesem Verfahren ergeben sich aber wiederum andere Nachteile. Je höher das beobachtete Gestirn steht, desto stärker wirken sich kleine Fehler bei der Justierung des Instruments aus. Diese ist an sich schon sehr heikel und wird leicht wieder etwas gestört, wenn sich das Instrument nicht sehr fest aufstellen lässt. Ein Fehler in der Zeitbestimmung wirkt sich hier ungünstiger auf das Azimut aus als bei geringen Höhen. Bei der Beobachtung der Sonne ist man in unseren oder weiter südlich gelegenen Breiten auf die Wintermonate beschränkt, sofern sich das Fernrohr nur bis zu $40^{\circ}$ neigen lässt. Fixsterne erster Größe lassen sich mit diesen Instrumenten zwar auch bei Tage beobachten. Muss man aber die Zeit aus korrespondierenden Sonnenhöhen bestimmen, so fällt auch noch ein möglicher Fehler der Sonnentafeln ins Gewicht. Hat man dagegen ein gutes Passageinstrument, so erübrigen sich in den meisten Fällen die Azimutbeobachtungen ganz, weil man mit diesem Instrument leicht und mit großer Genauigkeit einen künstlichen Gegenstand in die Meridianebene bringen kann. Aus diesen Gründen würden wir Azimutbestimmungen außerhalb der Kulmination und bei mäßigen Höhen vorziehen. Besteht noch eine kleine Unsicherheit in der Polhöhe oder in der Deklination des Gestirns, so kann man sie leicht unschädlich machen, indem man Beobachtungen beiderseits des Meridians in nahezu gleichen Abständen von ihm verbindet. Ebenso wenig fällt hier die zuvor erwähnte Ungleichförmigkeit der Azimutänderungen ins Gewicht, da sie sich leicht und mit größter Genauigkeit in Rechnung stellen lässt. In der Tat ist klar, dass man das Azimut zum Stundenwinkel $t+\theta$ nach dem Lehrsatz von \textsc{Taylor} in eine Reihe $A+B\theta+C\theta\theta+D\theta^3+\text{etc.}$ entwickeln kann, von der wenige Glieder genügen, sofern $\theta$ innerhalb enger Grenzen bleibt. Sind also insgesamt $n$ Beobachtungen gemacht worden, denen die Stundenwinkel $t+\theta$, $t+\theta'$, $t+\theta''$, $t+\theta'''$ usw. entsprechen, so ist das Mittel aller Azimute \SourceLeaf{593}offenbar
\[
\begin{aligned}&= A + \frac{\theta+\theta'+\theta''+\theta'''+\text{etc.}}{n}\cdot B\\[3pt]
&+ \frac{\theta\theta+\theta'\theta'+\theta''\theta''+\theta'''\theta'''+\text{etc.}}{n}\cdot C\\[3pt]
&+ \text{etc.}\end{aligned}
\]
Da man den zugrunde gelegten Wert von $t$ hier frei wählen kann, nimmt man am einfachsten denjenigen, der dem Mittel aller Beobachtungszeiten entspricht. Dadurch ergibt sich $\theta+\theta'+\theta''+\theta'''+\text{etc.}=0$, und folglich enthält das Mittel aller Azimute die Form.
\[
A + \frac{1}{n}\,C\,\Sigma\theta\theta + \frac{1}{n}\,D\,\Sigma\theta^3
\]
Denn in der Praxis wird man nie über die Glieder dritter Ordnung hinausgehen müssen. In weitaus den meisten Fällen genügt sogar die Korrektur zweiter Ordnung $\frac{1}{n}\,C\,\Sigma\theta\theta$, an deren Stelle man offenbar auch
\[
2\,C \times \frac{\Sigma\,2\sin\frac{1}{2}\theta^2}{n}
\]
nehmen darf, da der Fehler nur von vierter Ordnung ist. Dadurch kann man die bekannte Hilfstafel zur Reduktion von Zirkummeridianhöhen verwenden, die $2\sin\frac{1}{2}\theta^2$ unmittelbar in Sekunden angibt. Dass man beim dritten Glied, mit $D$ so, wie es sich aus der Entwicklung nach \textsc{Taylors} Lehrsatz ergibt, noch den Faktor $\frac{15^3}{206265^2}$ oder $\frac{900^3}{206265^2}$ oder $\frac{9000^3}{206265^2}$ hinzufügen muss, wenn man $\theta$ in Zeitsekunden, Zeitminuten oder Zehnern von Zeitminuten ausdrückt, versteht sich von selbst.

Dieses hier kurz beschriebene Verfahren stellt Herr \textsc{Soldner} in der vorliegenden Abhandlung vor. Es ist so einfach und liegt so nahe, dass man sich wundern muss, wie es mehreren praktischen Astronomen bei gleichem oder ganz ähnlichem Anlass entgehen konnte. Gerade deshalb verdient Herr \textsc{Soldner} für die Bekanntmachung des Verfahrens den Dank der praktischen Astronomen. Dies umso mehr, als die Entwicklungen geschickt und elegant durchgeführt und die Endergebnisse, das heißt die Ausdrücke für die Koeffizienten $B$, $C$, $D$ (denn so weit hat Herr \textsc{Soldner} die Entwicklung geführt), in eine sehr handliche Form gebracht sind.

\begingroup\rightskip=0pt plus 3em
Der Rezensent hat seit ein paar Jahren bei häufigen Azimutbeobachtungen mit einem \textsc{Reichenbach}'schn Theodoliten genau dieses Verfahren angewandt, allerdings ohne das in der Praxis meist überflüssige Glied dritter Ordnung zu berücksichtigen. Er hofft, dass es den astronomischen Lesern willkommen ist, hier noch einige Bemerkungen zu diesem Gegenstand zu finden. Er benutzte dabei die Formel
\[
2\,C = -\,\frac{\cos\varphi\cdot \operatorname{tang}\,\delta\cdot \sin A^2}{\sin t}
\left(1 + \frac{2\cos\varphi\cdot \sin 2A}{\sin t\cdot \sin 2\delta}\right)
\]
wobei $\varphi$, $\delta$, $t$, $A$ für Polhöhe, Deklination, Stnndenwinkel und Azimut stehen. Man kann die Formel auch so darstellen.
\[
a = \frac{\cos\varphi\cdot \operatorname{tang}\,\delta\cdot \sin A^2}{\sin t}, \qquad
2\,C = -\,a\left(1 + \frac{2a}{\operatorname{tang}\,A\cdot \sin\delta^2}\right)
\]
Diese Formel ist im Grunde identisch mit einer, die auch Herr \textsc{Soldner} auf S. 12 anführt und dort ebenfalls nur für den Fall empfiehlt, dass man das Glied dritter Ordnung nicht berücksichtigen will. Dort wird nur nicht erwähnt, dass es \emph{unter dieser Voraussetzung} bequemer ist, zur Berechnung des Azimuts selbst anstelle der \textsc{Neper}'schen Formeln die folgenden, an sich bekannten Formeln zu verwenden.
\SourceLeaf{594}
\[
\begin{gathered}
\frac{\operatorname{tang}\delta}{\cos t}=\operatorname{tang}\psi\\[5pt]
\frac{\cos\psi\cdot\operatorname{tang}t}{\sin(\varphi-\psi)}=\operatorname{tang}A
\end{gathered}
\]
Soll dagegen das Glied dritter Ordnung mit berücksichtigt werden, so benutzt man für $C$ und $D$ die Formeln von Herrn \textsc{Soldner} auf S. 9. An ihrer Form kann man auch noch eine kleine Änderung vornehmen, auf die wir weiter unten zu sprechen kommen.\par\endgroup

Wann immer man sich mit der Korrektur zweiter Ordnung begnügen kann, lässt sich sogar die Berechnung des Koeffizienten $C$ durch einen kleinen Kunstgriff ganz umgehen. Man kann ihn auch bei mancherlei anderen Gelegenheiten mit Vorteil anwenden, und wir wollen ihn hier noch anführen. Wie vorhin sei $t$ der Stundenwinkel zum Mittel aller Beobachtungszeiten, oder $\theta+\theta'+\theta''+\theta'''+\text{etc.}=0$, und
\[
\tau=\BareRoot{\frac{\theta\theta+\theta'\theta'+\theta''\theta''+\theta'''\theta'''+\text{etc.}}{n}}
\]
Das Mittel aller $n$ Azimute ist dann dem Mittel aus den beiden Azimuten gleichzusetzen, die den Stundenwinkeln $t-\tau$ und $t+\tau$ entsprechen. Man kann diese beiden Azimute ungefähr ebenso leicht direkt berechnen wie das eine Azimut $A$ und den Koeffizienten $C$. Um $\tau$ leichter zu bestimmen, darf man ohne Bedenken die bereits erwähnte Reduktionstafel benutzen. Zu dem in Zeit ausgedrückten $\tau$ gehört darin das Mittel der einzelnen $n$ Größen, die in derselben Tafel den in Zeit ausgedrückten Argumenten $\theta$, $\theta'$, $\theta''$, $\theta'''$ usw. entsprechen. Das Verfahren empfiehlt sich auch deshalb, weil man dabei ohne zusätzlichen Aufwand auf allerlei Begleitumstände Rücksicht nehmen kann, etwa bei der vorliegenden Aufgabe auf die Änderung der Deklination der Sonne, wenn es um die Sonne geht. Unmittelbar ist dieser Kunstgriff allerdings nur anwendbar, wenn die Stundenwinkel bekannt sind und die Azimute gesucht werden. Will man beobachtete Azimutunterschiede zur Zeitbestimmung nutzen, braucht man weitere Hilfsmittel, die hier nicht zu erläutern sind.

Betrachtet man die von Herrn \textsc{Soldner} behandelte Untersuchung aus rein mathematischer Sicht, so liefert sie im Grunde die Lösung der folgenden Aufgabe. Sind in einem sphärischen Dreieck zwei Seiten unveränderlich, so sind die endlichen Änderungen eines anliegenden Winkels als Reihe nach aufeinanderfolgenden Potenzen der endlichen Änderungen des eingeschlossenen Winkels auszudrücken. Das Gegenstück wäre, unter denselben Bedingungen die endlichen Änderungen der dritten Seite auf ähnliche Weise darzustellen. Auch diese Aufgabe muss man bei mehreren Gelegenheiten lösen, um Beobachtungen mit Vervielfältigungsinstrumenten genau zu reduzieren. Mit diesen beiden Aufgaben und verschiedenen Anwendungen befasst sich die Schrift Nr. II. Ihr Verfasser kannte die Abhandlung von Herrn \textsc{Soldner} bereits, wie man aus mehreren zum Teil sehr unpassenden Äußerungen darüber sieht. Dazu gehört auf S. 29 die seltsame Behauptung, das Verfahren von \textsc{Soldner} sei nichts weiter als die von \textsc{Svanberg} angewandte Methode. Wer beide Methoden vergleicht, erkennt ohne Weiteres, dass dieses Urteil jeder Grundlage entbehrt. Wir brauchen daher nicht länger dabei zu verweilen. Was die Ausführung angeht, können wir der Schrift von \textsc{Stefenelli} nicht dasselbe Lob erteilen wie der von \textsc{Soldner}. Die beiden Reihen, mit denen Herr \textsc{Stefenelli} die beiden Aufgaben zu lösen glaubt, sind beide schon in den Gliedern dritter Ordnung \emph{unrichtig}. Wenn wir uns recht erinnern, wurden sie schon vor einiger Zeit in einem gelehrten Blatt unter dem unpassenden Titel neuer Entdeckungen angekündigt. Dass diese Glieder in der Praxis meist unbedeutend sind, entschuldigt den Verfasser nicht. Er hätte sich auf die Glieder zweiter Ordnung beschränken können. Wenn er es aber einmal unternahm, darüber \SourceLeaf{595}hinauszugehen, muss man richtige Ergebnisse verlangen. Dies umso mehr, als Herr \textsc{Soldner} die eine Aufgabe bereits bis zur dritten Ordnung vollständig gelöst hatte und überhaupt das ganze Verdienst dieser an sich gar nicht schwierigen Entwicklungen allein in ihrer Richtigkeit und Handlichkeit besteht. Das Einzige, was wir zum Lob der Schrift von \textsc{Stefenelli} anführen können, ist, dass dort das noch richtige Glied zweiter Ordnung (bei der ersten Aufgabe, die bei Herrn \textsc{Stefenelli} die zweite ist) eine etwas bequemere Form hat als bei Herrn \textsc{Soldner}. Herr \textsc{Stefenelli} muss aber übersehen haben, dass sich der Reihe von \textsc{Soldner} durch die einfachen, auf bekannten Gleichungen beruhenden Substitutionen
\[
\begin{gathered}
\sin\varphi+\sin\delta=\frac{2\cos\frac{1}{2}z^2\cdot\sin2\gamma}{\sin t}\\[5pt]
\sin\varphi-\sin\delta=\frac{2\sin\frac{1}{2}z^2\cdot\sin2\beta}{\sin t}
\end{gathered}
\]
sofort dieselbe Form geben lässt. Hätte er das bemerkt, so hätte er durch den Vergleich des Glieds dritter Ordnung erkennen können, dass sein eigenes unrichtig ist. Man begreift kaum, wie er übersehen konnte, dass die Verwechslung von Bögen mit ihren Sinus und Tangenten auf S.\,3 notwendig schon Fehler dritter Ordnung hervorbringen musste.

\begingroup\rightskip=0pt plus 3em
Einer Schlussbemerkung zufolge hat Herr \textsc{Stefenelli} einen Teil seiner Schrift der mathematisch-physikalischen Klasse der Münchner Akademie vorgelegt, anscheinend sogar in der Erwartung, sie in die Denkschriften der Akademie aufgenommen zu sehen. Das ist freilich nicht geschehen. Ob die Akademie ein Urteil darüber abgegeben hat, wird nicht gesagt. Wir wären bei der Besprechung weniger streng gewesen oder hätten die Schrift lieber ganz übergangen, da sie, soweit wir wissen, der erste Versuch des Verfassers zu sein scheint, wenn sein unüberlegter Angriff auf die Abhandlung von \textsc{Soldner} keine ernstliche Rüge verdient hätte. Um des Verfassers selbst willen ist es zu bedauern, dass er nicht dazu veranlasst wurde, seinen Versuch erst weiter ausreifen zu lassen, ehe er damit an die Öffentlichkeit trat.\par\endgroup
